\section{Diseño criptográfico del prototipo}\label{sec:diseno-cripto}

Las decisiones de esta sección aplican los estándares del NIST, las RFC del IETF y las guías de OWASP del marco teórico a las restricciones jurídicas y técnicas del prototipo. Cubren cuatro servicios: integridad y autoría mediante firma digital, existencia en el tiempo mediante sellado, confidencialidad en reposo mediante cifrado, y autenticación con control de acceso.

\subsection{Nivel de seguridad objetivo}\label{sec:nivel-seguridad}

El prototipo adopta un nivel de fortaleza de 128 bits conforme a NIST SP 800-57 \cite{barker2020}. Ninguna primitiva ofrece menos, de modo que no existe un eslabón débil en el que un atacante racional pueda concentrar su esfuerzo. Es el primer nivel que NIST acepta para uso general sin restricción temporal; el de 112 bits, que corresponde a RSA-2048, queda como uso heredado y no se recomienda para proteger información nueva después de 2030. La diferencia importa porque un documento firmado puede presentarse como prueba judicial años después. La e.firma del SAT usa RSA-2048; la CA interna permite elegir parámetros por encima de ese piso. La \cref{tab:cripto-primitivas} resume cada primitiva. AES-256 supera el objetivo; se eligió por su margen adicional, con un costo de rendimiento pequeño.

\begin{table}[!htbp]
  \centering
  \caption{Primitivas criptográficas, uso, parámetros, fuerza y norma.}
  \label{tab:cripto-primitivas}
  \footnotesize
  \begin{tabularx}{\linewidth}{@{}L{2.55cm} X L{3.3cm} L{1.75cm} L{2.0cm}@{}}
    \toprule
    \textbf{Primitiva} & \textbf{Uso en el sistema} & \textbf{Parámetros} & \textbf{Fuerza (bits)} & \textbf{Norma} \\
    \midrule
    SHA-256 & Huella de cada versión, de las firmas y de los consentimientos; digest que sella la TSA & Salida de 32 B & 128 (colisión) & FIPS 180-4 \cite{nist2015fips1804} \\
    RSA-3072, PKCS\#1 v1.5 con SHA-256 & Firma del PDF y del consentimiento, certificados, CRL, tokens de la TSA y acceso del Owner por certificado & Clave de 3072 bits; firma de 384 B & 128 & SP 800-57 \cite{barker2020} \\
    X.509 v3 y CRL & Certificados de la CA interna y su revocación & Raíz de 5 años, finales de 1 año, CRL obligatoria & 128 & RFC 5280 \cite{cooper2008} \\
    Sello RFC 3161 & Existencia del digest del PDF en un instante & SHA-256; firma RSA-3072 de la TSA & 128 & RFC 3161 \cite{adams2001} \\
    AES-256-GCM & Cifrado de cada versión documental, envoltura de las DEK y secretos TOTP & Clave de 32 B, nonce aleatorio de 12 B, etiqueta de 16 B & 256 & SP 800-38D \cite{dworkin2007} \\
    Envoltura KEK/DEK & Separar la clave de cada versión de la clave maestra; rotar sin recifrar & Una DEK de 32 B por versión & 256 & SP 800-38F \cite{dworkin2012} \\
    Argon2id & Contraseñas y códigos de recuperación & $m = 256$ MiB, $t = 3$, $p = 1$; sal de 16 B; salida de 32 B; resistente a memoria & Depende de la contraseña & RFC 9106 \cite{biryukov2021} \\
    TOTP con HMAC-SHA1 & Segundo factor de acceso & 6 dígitos, periodo de 30~s, secreto de 20 B & No aplica & RFC 6238 \cite{mraihi2011}; SP 800-107 \cite{dang2012} \\
    \bottomrule
  \end{tabularx}
\end{table}

\subsection{Qué recibe cada parte}\label{sec:cripto-fronteras}

Las zonas y sus cruces se describen en la \cref{fig:zonas-confianza}; aquí se precisa qué material criptográfico ve cada parte. El servidor de la API forma parte de la frontera de confidencialidad: recibe los documentos en claro, los cifra con claves que solo él maneja y los descifra para validar su formato y entregarlos. La protección del lado del navegador se limita a la clave de firma del titular, que nunca sale de su equipo. La \cref{tab:cripto-visibilidad} lo resume.

\begin{table}[!htbp]
  \centering
  \caption{Qué recibe o conserva cada parte del sistema y qué queda fuera de su alcance.}
  \label{tab:cripto-visibilidad}
  \footnotesize
  \begin{tabularx}{\linewidth}{@{}L{2.4cm} X L{4.6cm}@{}}
    \toprule
    \textbf{Parte} & \textbf{Recibe o conserva} & \textbf{No recibe ni conserva} \\
    \midrule
    Navegador del titular & Clave privada cifrada y su contraseña durante la operación local; bytes preparados por el servidor (PDF, declaraciones, desafíos); firmas resultantes & No guarda la clave ni su contraseña en borradores, almacenamiento local ni registros \\
    API (servidor) & Contraseña y código TOTP al iniciar sesión; documentos en claro al cargarlos y entregarlos; KEK; certificados públicos, firmas y consentimientos & Claves privadas de los titulares, sus contraseñas y contenedores PKCS\#12 \\
    PostgreSQL & Documentos cifrados y DEK envueltas; hashes Argon2id; secretos TOTP cifrados; certificados, firmas, consentimientos, tokens RFC 3161, revisiones de confianza y cadena de auditoría & KEK, documentos en claro, contraseñas y claves privadas de los titulares \\
    Redis & Sesiones opacas y revocables, desafíos MFA, contadores de intentos y huellas de códigos TOTP aceptados & El código TOTP en claro; solo su huella \\
    CA interna & CSR público con nombre, organización y país declarados; certificados emitidos; revocaciones y CRL & Clave privada del titular y su contraseña \\
    TSA local & Digest SHA-256 del documento & El documento y el consentimiento \\
    Proveedor de correo & Dirección del destinatario y un mensaje genérico; en invitaciones y recuperación de contraseña, además, el enlace con su token & Datos del expediente, documentos y adjuntos \\
    \bottomrule
  \end{tabularx}
\end{table}

\subsection{Firma digital}\label{sec:firma-digital}

El esquema de firma emplea RSA-3072 con SHA-256 y relleno PKCS\#1 v1.5. Se descartó ECDSA P-256 porque la Política de Certificados de la ACR2-SE \cite{economia_acr2se} exige claves RSA con ese algoritmo en la jerarquía PKI mexicana; ECDSA no está contemplado para certificados legalmente reconocidos bajo un PSC mexicano. Por la misma razón se prefirió PKCS\#1 v1.5 sobre PSS: coincide con la ACR2-SE y la e.firma, lo que permitiría migrar a certificados de un PSC sin rediseñar el verificador.

La autoridad certificadora interna usa OpenSSL y una jerarquía sin CA intermedia: la raíz emite directamente los certificados finales (\cref{fig:cripto-pki}). Para la firma individual, el titular genera su clave RSA-3072 y la solicitud pública de certificado (CSR) en su equipo, con OpenSSL o con el ayudante local del navegador. La CA comprueba la solicitud y emite un certificado exclusivamente de firma, sin generar ni recibir la clave privada. El portal administrativo separado (\cref{sec:impl-portal-ca}) exige además la revisión y confirmación del operador.

Los certificados tienen perfiles separados: uno de firma documental y otro con el que el Owner puede acceder al sistema, y el verificador de firmas solo admite el primero. Cambiar de emisor requiere adaptar y validar la política de confianza. Al vencer un certificado se emite otro y, cuando corresponde, se registra un nuevo vínculo, sin modificar la evidencia histórica.

\begin{figure}[!htbp]
  \centering
  \begin{tikzpicture}[
  hoja/.style={caja, text width=3.3cm, minimum height=0.95cm},
  bloque/.style={caja, text width=3.4cm, minimum height=0.85cm},
  uso/.style={cajaClara, text width=2.9cm, minimum height=0.85cm}]

  % Zonas
  \draw[zona] (0,0) rectangle (16,-3.4);
  \draw[zona] (0,-3.9) rectangle (7.75,-7.85);
  \draw[zona] (8.25,-3.9) rectangle (16,-7.85);
  \node[zonaLabel] at (0,0) {CA INTERNA (PORTAL SEPARADO CON MFA)};
  \node[zonaLabel, anchor=south west] at (0,-7.85) {EQUIPO DEL TITULAR};
  \node[zonaLabel, anchor=south east] at (16,-7.85) {SERVIDOR DEL DESPACHO};

  % Jerarquía de la CA
  \node[cajaOsc, text width=5.0cm, minimum height=0.8cm] (raiz) at (8,-0.85)
    {Raíz interna autofirmada\\ RSA-3072, vigencia de 5 años};
  \node[hoja] (c2) at (2.15,-2.6) {Acceso del Owner\\ desde CSR (1 año)};
  \node[hoja] (c1) at (6.05,-2.6) {Firma documental\\ desde CSR (1 año)};
  \node[hoja] (c3) at (9.95,-2.6) {TSA local\\ solo sellado (1 año)};
  \node[cajaClara, text width=3.3cm, minimum height=0.95cm] (c4) at (13.85,-2.6)
    {CRL firmada por la raíz\\ (vigencia de 7 días)};
  \draw[ESCOMAzulOscuro, line width=0.6pt] (raiz.south) -- (8,-1.65);
  \draw[ESCOMAzulOscuro, line width=0.6pt] (2.15,-1.65) -- (13.85,-1.65);
  \foreach \c/\x in {c1/6.05,c2/2.15,c3/9.95,c4/13.85}
    \draw[flecha] (\x,-1.65) -- (\c.north);

  % Titular: un uso por certificado
  \node[uso] (acc) at (2.15,-5.0) {Firma desafíos de acceso del Owner};
  \node[uso] (fir) at (6.05,-5.0) {Firma vínculo, PDF y consentimiento};
  \node[cajaOsc, text width=3.3cm, minimum height=0.85cm] (priv) at (2.15,-6.6)
    {Clave RSA-3072 generada con el ayudante local u OpenSSL; cifrada, nunca sale del equipo};
  \draw[flecha] (priv.north) -- (acc.south);
  \draw[flecha] (priv.north east) -- (fir.south west);

  % Intercambio con la CA: CSR de subida, certificado de bajada
  \foreach \c/\b/\x in {c2/acc/2.15,c1/fir/6.05} {
    \draw[flecha] (\x-0.75,-4.575) -- node[etq, pos=0.6]{CSR} (\x-0.75,-3.075);
    \draw[flecha] (\x+0.75,-3.075) -- node[etq, pos=0.4]{certificado} (\x+0.75,-4.575);
  }

  % Servidor
  \node[bloque] (tsa) at (9.95,-4.95) {TSA local (OpenSSL)\\ sella digests SHA-256};
  \node[bloque] (conf) at (13.85,-4.95) {Revisión de confianza\\ (raíz + CRL publicadas)};
  \node[bloque] (vin) at (13.85,-6.95) {Vínculo cuenta y certificado;\\ verifica las firmas};
  \draw[flecha] (conf) -- node[etq, right]{se captura} (vin);
  \draw[flecha] (c3) -- node[etq, pos=0.45]{certificado} (tsa);
  \draw[flecha] (c4) -- node[etq, pos=0.45]{publicación} (conf);

  % Material público hacia el servidor
  \draw[flecha] ($(fir.south)+(0.6,0)$) |- node[etq, below, pos=0.8, align=center]
    {certificado y firmas\\ (material público)} (vin.west);
\end{tikzpicture}

  \caption{PKI interna. La raíz emite los certificados finales y publica la CRL; los certificados de firma y de acceso nacen de un CSR generado en el equipo del titular, de modo que la CA nunca recibe la clave privada. Solo el certificado de firma documental firma vínculo, PDF y consentimiento.}
  \label{fig:cripto-pki}
\end{figure}

La revocación se comprueba con una CRL obligatoria, firmada por la raíz y vigente en el instante evaluado, conforme a RFC 5280 \cite{cooper2008}. La raíz y la CRL se publican en el servidor como una revisión explícita de confianza. Un nuevo vínculo o consentimiento captura esa revisión y la vuelve a comparar al confirmar. Las revisiones históricas conservan su material, de modo que las pruebas originales se evalúan sin sustituirlas por la CRL actual.

Una firma con certificado de la CA interna usa los mismos mecanismos técnicos que una con certificado de la e.firma o de un PSC: RSA-3072, SHA-256, X.509 v3 y verificación de cadena. Cambia el emisor y, con él, el nivel de confianza jurídica. La CA interna permite demostrar integridad y autoría en el alcance académico, pero no se presenta como equivalente a una firma avanzada.

La custodia del titular se separa de la operación de la CA. El navegador permite usar una clave cifrada y su contraseña dentro de un proceso local aislado, o importar firmas hechas con una herramienta externa. La API recibe solamente material público. El vínculo explícito identifica una cuenta del sistema y exige prueba de posesión de la clave; el nombre o el correo del certificado no crean ese vínculo ni acreditan identidad civil. La \cref{sec:impl-criptografia-local} detalla el formato cifrado y la confianza depositada en el código del navegador.

La autorización de firma se modela aparte de la posesión de la clave. Primero, una declaración canónica vincula el certificado con la cuenta y con una revisión de confianza. Después, cada operación documental requiere dos firmas del titular (\cref{fig:cripto-consentimiento}): la del PDF y la de un consentimiento que fija expediente, documento, versión, digest del PDF y digest de esa primera firma. Por ello, un consentimiento no puede trasladarse a otra cuenta ni a otra versión, aunque los bytes del PDF coincidan. Ninguna de las dos firmas confirma por sí sola la operación.

\begin{figure}[!htbp]
  \centering
  \begin{tikzpicture}[
  dato/.style={caja, text width=2.7cm, minimum height=0.85cm},
  firma/.style={cajaOsc, text width=2.7cm, minimum height=0.85cm},
  oper/.style={cajaClara, text width=3.0cm, minimum height=0.85cm},
  campo/.append style={font=\diag}]

  % Consentimiento: contexto común con el vínculo
  \node[anchor=south west, font=\diag\bfseries, text=ESCOMAzulOscuro, inner sep=1pt]
    at (0,0.95) {Consentimiento (258 B). Parte 1: cuenta y confianza, con la misma estructura que la declaración de vínculo (142 B)};
  \Tira{0}{0}{2.3/campo/{encabezado: versión y propósito}//,
              1.7/campo/{despliegue}//,
              1.9/campoRel/{huella de la raíz}//,
              1.9/campoRel/{revisión de confianza}//,
              1.4/campoRel/{cuenta}//,
              2.4/campo/{revisión de cuenta y generación de acceso}//,
              1.5/campoRel/{vínculo}//,
              2.6/campoRel/{huella del certificado}//}

  % Consentimiento: datos de la operación
  \node[anchor=south west, font=\diag\bfseries, text=ESCOMAzulOscuro, inner sep=1pt]
    at (0,-0.75) {Parte 2: operación documental (116 B)};
  \Tira{0}{-1.65}{2.2/campoRel/{operación}//,
              2.2/campoRel/{expediente}//,
              2.2/campoRel/{documento}//,
              1.6/campoRel/{versión}//,
              3.6/campoFirma/{digest SHA-256 del PDF}//,
              3.9/campoFirma/{digest SHA-256 de la firma del PDF}//}

  % Flujo de firmas en el equipo del titular
  \node[dato] (pdf) at (9.75,-3.75) {PDF de la versión exacta};
  \node[firma, text width=3.1cm] (fpdf) at (13.95,-3.75) {Firma 1: firma separada del PDF (384 B)};
  \draw[flecha] (pdf) -- node[etq, above]{firma 1} (fpdf);
  \draw[flecha] (9.2,-3.325) -- node[etq, left]{SHA-256} (9.2,-1.65);
  \draw[flecha] (13.1,-3.325) -- node[etq, left]{SHA-256} (13.1,-1.65);

  \node[oper, text width=3.4cm] (f2) at (2.2,-3.75) {El titular firma los 258~B (partes 1 y 2) con la misma clave};
  \node[firma] (fcon) at (6.2,-3.75) {Firma 2: firma del consentimiento (384 B)};
  % Llave que abarca las dos partes
  \draw[ESCOMAzulOscuro, line width=0.6pt] (-0.05,-1.65) -- (-0.3,-1.65);
  \draw[flecha] (-0.05,0.85) -- (-0.3,0.85) |- (f2.west);
  \draw[flecha] (f2) -- (fcon);

  % Sello de tiempo
  \node[caja, text width=5.4cm, minimum height=0.85cm] (tsa) at (9.75,-5.35)
    {TSA local: token RFC 3161 sobre el digest del PDF; no cubre el consentimiento};
  \draw[flecha] (pdf.south) -- node[etq, right]{solo el digest} (tsa.north);
  \node[cajaGris, text width=6.3cm, minimum height=0.85cm, align=left] at (3.5,-5.35)
    {El servidor recibe el consentimiento y las dos firmas, nunca la clave. Reconstruye el contexto vigente y exige igualdad exacta.};
\end{tikzpicture}

  \caption{Qué se firma en la firma individual. La firma del PDF entra al consentimiento a través de su digest, y la misma clave firma ambos mensajes en el equipo del titular. El sello RFC 3161 cubre solo el digest del PDF.}
  \label{fig:cripto-consentimiento}
\end{figure}

\subsection{Sellado de tiempo}\label{sec:sellado-tiempo}

El sellado de tiempo se aísla en un componente intercambiable. El servidor usa una TSA local con OpenSSL; un adaptador para un PSC externo existe, pero no está seleccionado. En ambos casos el flujo sigue RFC 3161 \cite{adams2001}: se calcula el SHA-256 del documento firmado, se solicita el sello al emisor configurado y la respuesta se guarda como archivo independiente (.tsr) en formato ASN.1 DER. Así, la TSA local puede sustituirse por Cincel, SeguriData, CECOBAN u otra TSA compatible sin modificar la lógica central. La equivalencia es de protocolo, no de confianza: solo un PSC autorizado y contratado produce la constancia prevista por la NOM-151-SCFI-2016.

Cincel, PSC acreditado con entorno de pruebas RFC 3161 \cite{cincel2024}, no se usa como evidencia del prototipo y queda como adaptador opcional por lo expuesto en la \cref{sec:riesgo-psc-materializado}.

Guardar el token como archivo .tsr, en lugar de incrustarlo en CAdES-T o PAdES-T, permite inspeccionar por separado documento, firma, sello y certificado con herramientas estándar. CAdES-T y PAdES-T exigirían ingeniería de formatos (estructuras CMS o firmas internas del PDF) sin aportar al objetivo criptográfico. El alcance jurídico depende de la autoridad emisora y su cadena de confianza \cite{nom151_2016}, no del contenedor; un despliegue productivo podría migrar a CAdES-T o PAdES-T reutilizando el mismo token.

El paquete de evidencia es un ZIP con el PDF original, su firma separada, el token RFC 3161 y la confianza capturada. En una firma individual agrega las declaraciones de vínculo y consentimiento con sus firmas, la raíz y la CRL del registro y una representación pública de sus tiempos. Las firmas individuales nuevas añaden un manifiesto sellado, la lista de revocación del emisor de la TSA y un recibo de aceptación (\cref{fig:cripto-manifiesto}). Sus instrucciones permiten comprobar todo con OpenSSL sin acceso a la plataforma, siguiendo los cuatro pasos de la \cref{fig:cripto-sello}; el mismo paquete también puede validarse en el navegador sin enviarlo al servidor (\cref{sec:validacion-navegador}). El sello del servidor nunca sustituye la firma del titular.

\begin{figure}[!htbp]
  \centering
  \begin{tikzpicture}[
  ent/.style={caja, text width=3.2cm, minimum height=0.7cm},
  pasoV/.style={cajaClara, text width=3.2cm, minimum height=1.45cm}]

  % Entradas
  \node[ent] (tsr) at (1.75,0) {Respuesta del sello (.tsr)};
  \node[ent] (pdf) at (6.0,0) {Documento original};
  \node[ent] (anc) at (14.25,0) {Ancla de confianza};

  % Pasos
  \node[pasoV] (p1) at (1.75,-1.75) {Extraer el token de la respuesta};
  \node[pasoV] (p2) at (6.0,-1.75) {Recalcular SHA-256 del documento y compararlo con la huella del token};
  \node[pasoV] (p3) at (10.0,-1.75) {Verificar la firma de la TSA sobre el token};
  \node[pasoV] (p4) at (14.25,-1.75) {Validar el certificado de la TSA contra el ancla};
  \foreach \p/\n in {p1/1,p2/2,p3/3,p4/4}
    \node[numero] at (\p.north west) {\n};
  \draw[flecha] (tsr) -- (p1);
  \draw[flecha] (pdf) -- (p2);
  \draw[flecha] (anc) -- (p4);
  \draw[flecha] (p1) -- (p2);
  \draw[flecha] (p2) -- (p3);
  \draw[flecha] (p3) -- (p4);

  % Resultados
  \node[cajaError, text width=11.6cm, minimum height=0.7cm] (rech) at (5.95,-3.55)
    {Si cualquier paso falla, la verificación del sello se rechaza};
  \foreach \p in {p1,p2,p3} \draw[flechaError] (\p.south) -- (\p.south |- rech.north);
  \draw[flechaError] ($(p4.south)+(-1.3,0)$) |- (rech.east);
  \node[cajaOsc, text width=2.3cm, minimum height=0.7cm] (ok) at (14.8,-3.55)
    {Instante válido para ese digest};
  \draw[flecha] ($(p4.south)+(0.55,0)$) -- (ok.north -| {$(p4.south)+(0.55,0)$});

  % Anclas posibles
  \node[cajaGris, text width=15.4cm, align=left, anchor=north west] at (0,-4.45)
    {\textbf{Ancla según el emisor.} TSA local: la raíz interna del prototipo. Constancia externa: la cadena del PSC hasta la ACR2-SE.};
\end{tikzpicture}

  \caption{Verificación de un sello RFC 3161 en cuatro pasos. La huella se compara antes de revisar la firma y la cadena; el ancla depende de quién emitió el sello.}
  \label{fig:cripto-sello}
\end{figure}

El instante del token corresponde al digest del PDF, no al consentimiento personal. La aceptación individual queda en un recibo y un evento auditado, bajo el reloj y la frontera de confianza de la plataforma. Por ello, un sello válido del contenido no demuestra por sí solo que el titular autorizó esa versión en ese instante. La verificación conserva las pruebas históricas y presenta aparte la vigencia y revocación actuales del certificado y el posible retiro del vínculo.

\subsection{Validación de evidencia en el navegador}\label{sec:validacion-navegador}

La ruta con OpenSSL exige una terminal, conocer los comandos y comparar a mano campos binarios. Para que cualquier persona con acceso al documento pueda revisar un paquete exportado, Qadra incorpora un validador en la pestaña \emph{Evidencia} de cada versión documental. La persona elige un ZIP guardado en su equipo y obtiene, en lenguaje llano, qué se comprobó, qué falló y qué no se pudo decidir. La ruta con OpenSSL se conserva sin cambios (\cref{anexo:verificacion}); el validador es una segunda forma de llegar a las mismas conclusiones, no un sustituto.

\paragraph{Por qué en el navegador.} El paquete nunca sale del equipo de la persona: el validador lo lee, lo revisa y muestra el resultado sin enviarlo al servidor (\cref{fig:cripto-validador-bloques}). Hay tres razones. Primera, el paquete es evidencia que ya está en manos de quien la revisa, y devolverlo al servidor que la produjo haría que el sistema juzgara su propia evidencia. Segunda, el servidor no tiene que recibir ni conservar archivos de terceros, de modo que la validación no abre una nueva vía de carga. Tercera, el resultado no depende de la respuesta del servidor: solo se le consulta qué raíces admite hoy y, si no responde, la validación continúa.

\begin{figure}[!htbp]
  \centering
  % Fuente: docs/adr/0085-in-browser-evidence-validation.md;
% web/apps/qadra/src/components/EvidenceValidation.svelte,
% web/apps/qadra/src/lib/evidence/, web/apps/qadra/src/lib/evidence-verify-*.mjs;
% crates/evidence-verify, crates/evidence-verify-wasm,
% crates/web/src/evidence_trust/mod.rs y crates/application/src/document_signing/.
\begin{tikzpicture}[
  a/.style={caja, text width=2.6cm, minimum height=1.05cm},
  b/.style={caja, text width=3.3cm, minimum height=1.05cm},
  s/.style={caja, text width=3.6cm, minimum height=1.05cm}]

  % Zonas
  \draw[zona] (0,0) rectangle (10.6,-7.6);
  \draw[zona, fill=white] (3.85,-4.2) rectangle (10.4,-7.35);
  \draw[zona] (11.5,0) rectangle (16,-7.6);
  \node[zonaLabel] at (0,0) {NAVEGADOR · PÁGINA DE QADRA};
  \node[zonaLabel, anchor=north east] at (10.4,-4.2) {WEB WORKER · UNO POR VALIDACIÓN};
  \node[zonaLabel] at (11.5,0) {SERVIDOR};

  % Navegador
  \node[a, cajaClara] (zip) at (1.55,-1.2) {\textbf{Paquete ZIP} guardado en el equipo};
  \node[b] (panel) at (5.75,-1.2) {\textbf{Pestaña Evidencia}: elige el ZIP, calcula su huella, arma la política y permite cancelar};
  \node[a, cajaOsc] (fijas) at (1.55,-3.0) {\textbf{Raíces fijas} de la versión de Qadra};
  \node[b] (cliente) at (5.75,-3.0) {\textbf{Cliente del worker}: crea un worker por validación y descarta respuestas tardías};
  \node[a, cajaClara] (inf) at (1.55,-5.4) {\textbf{Informe} global y por componente; descarga en JSON};

  % Worker
  \node[caja, text width=2.7cm, minimum height=1.05cm] (wasm) at (5.75,-5.4) {\textbf{Módulo WebAssembly}: recibe bytes y política; informa etapas};
  \node[cajaOsc, text width=2.3cm, minimum height=1.05cm] (motor) at (9.1,-5.4) {\textbf{Motor} en Rust, sin red, reloj ni azar};
  \node[font=\diag, text=ESCOMGris, text width=6.0cm, align=center] at (7.1,-6.85)
    {El worker se cierra al entregar el informe, al fallar o al cancelar};

  \draw[flecha] (zip) -- node[etq, above=2pt]{elige} (panel);
  \draw[flecha] (panel) -- node[etq, right=2pt]{bytes y política} (cliente);
  \draw[flecha] (fijas.east) -- (3.6,-3.0) |- node[etq, pos=0.25, right=1pt, align=left]{raíces\\ fijas} ([yshift=-2.5mm]panel.west);
  \draw[flecha] (cliente) -- node[etq, left=2pt]{inicia} (wasm);
  \draw[flechaD] (wasm) -- (motor);
  \draw[flecha] (wasm) -- node[etq, above=2pt]{informe} (inf);

  % Servidor
  \node[s] (ruta) at (13.75,-1.2) {\textbf{Ruta de confianza}: solo lectura, para cualquier sesión, sin expediente};
  \node[s] (exp) at (13.75,-2.95) {\textbf{Exportación}: arma el ZIP con los bytes guardados};
  \node[s] (acep) at (13.75,-4.65) {\textbf{Aceptación de la firma}: escribe el manifiesto con el codificador del motor y lo hace sellar por la TSA};
  \node[s, cajaOsc] (pg) at (13.75,-6.4) {\textbf{PostgreSQL}: recibos y manifiestos inmutables};
  \draw[flecha] (acep) -- (pg);
  \draw[flecha] (pg.east) -- ++(0.2,0) |- (exp.east);

  % Cruces
  \draw[flecha] ([yshift=1.5mm]panel.east) -- node[etq, above=1pt, pos=0.5]{raíces vigentes} ([yshift=1.5mm]ruta.west);
  \draw[flechaSuave] ([yshift=-1.5mm]ruta.west) -- node[etq, below=1pt, pos=0.5]{estado actual y notas de confianza} ([yshift=-1.5mm]panel.east);

  % Frontera
  \draw[IPNGuinda, dashed, line width=0.9pt] (11.05,0.15) -- (11.05,-7.75);
  \node[font=\diag\bfseries, text=IPNGuinda, rotate=90, fill=white, inner sep=1.5pt, align=center]
    at (11.05,-4.6) {el ZIP no se envía\\ al servidor};
\end{tikzpicture}

  \caption{Bloques del validador. El paquete se procesa dentro del navegador, en un worker que se crea para cada validación; el servidor solo informa qué raíces admite hoy, y esa respuesta nunca convierte una raíz en ancla.}
  \label{fig:cripto-validador-bloques}
\end{figure}

El motor de verificación está escrito en Rust y se compila a WebAssembly \cite{w3c_wasm2019}, el formato de código que los navegadores ejecutan de forma aislada. Se ejecuta en un Web Worker \cite{whatwg_workers}, un hilo separado de la página, para que una revisión pesada no congele la interfaz y para que pueda cancelarse cerrando ese hilo. Es un módulo distinto del ayudante que abre claves privadas (\cref{sec:impl-criptografia-local}): validar evidencia pública no necesita ninguna clave. El motor no usa red, reloj ni números aleatorios; recibe los bytes del paquete y una política explícita, que incluye el instante de evaluación, y con las mismas entradas produce siempre el mismo informe. El servidor reutiliza ese mismo motor para escribir el manifiesto y el recibo de la firma individual, así que escribe exactamente lo que el validador lee.

\paragraph{En qué raíces se confía.} Una raíz no se vuelve confiable por viajar dentro del ZIP, porque quien altera un paquete también puede cambiar su raíz. Cada versión de Qadra incluye una lista fija de huellas SHA-256 de raíces admitidas, cada una con los propósitos que puede cumplir: sello del servidor, firma individual o sello de tiempo. Una raíz del paquete solo se admite si su huella está en esa lista para el propósito requerido. Además, el panel consulta una ruta autenticada del servidor que devuelve las raíces y listas de revocación que admite en ese momento, y contrasta ambas fuentes (\cref{fig:cripto-validador-raices}). Solo la lista fija convierte una raíz en ancla y permite un resultado completo. Una raíz que la ruta devuelve sin estar en la lista queda como confianza sin confirmar, con resultado \emph{Parcial}; si ninguna de las dos fuentes la reconoce, el paquete no es compatible. Si la ruta no responde, solo cambia el estado actual. El validador nunca sigue direcciones encontradas dentro del paquete.

\begin{figure}[!htbp]
  \centering
  % Fuente: docs/adr/0085-in-browser-evidence-validation.md (raíces de confianza);
% web/apps/qadra/src/lib/evidence/pinned-roots.json, policy.mjs y notes.mjs;
% crates/web/src/evidence_trust/mod.rs.
\begin{tikzpicture}[
  q/.style={cajaClara, text width=3.3cm, minimum height=1.2cm},
  nota/.style={cajaGris, text width=4.6cm, minimum height=0.85cm, align=left}]

  \node[caja, text width=2.7cm, minimum height=1.2cm] (raiz) at (1.5,-1.0)
    {\textbf{Raíz que trae el paquete} o que aparece dentro de un sello};
  \node[q] (fija) at (5.75,-1.0)
    {¿Su huella SHA-256 está en la lista fija de esta versión para el propósito requerido?};
  \node[q, text width=2.8cm] (ruta) at (10.0,-1.0)
    {¿La ruta de confianza del servidor también la devuelve?};
  \node[cajaOsc, text width=2.3cm, minimum height=1.2cm] (ok) at (14.5,-1.0)
    {\textbf{Comprobada}\\ Raíz fija en Qadra};
  \node[caja, text width=2.8cm, minimum height=1.0cm] (dif) at (10.0,-3.05)
    {\textbf{Sigue comprobada}; el informe señala la diferencia};
  \node[q, text width=2.8cm] (ruta2) at (5.75,-3.15)
    {¿La ruta del servidor la devuelve?};
  \node[caja, text width=2.7cm, minimum height=1.0cm] (sin) at (1.5,-3.15)
    {\textbf{Confianza sin confirmar}: falta evidencia para decidir (Parcial)};
  \node[caja, text width=3.3cm, minimum height=1.0cm] (nc) at (5.75,-5.2)
    {\textbf{No compatible}, no Fallido: ninguna fuente reconoce la raíz};

  \draw[flecha] (raiz) -- (fija);
  \draw[flecha] (fija) -- node[etq, above=1pt]{sí} (ruta);
  \draw[flecha] (ruta) -- node[etq, above=1pt]{sí} (ok);
  \draw[flecha] (ruta) -- node[etq, right=1pt]{no} (dif);
  \draw[flecha] (fija) -- node[etq, right=1pt]{no} (ruta2);
  \draw[flecha] (ruta2) -- node[etq, above=1pt]{sí} (sin);
  \draw[flecha] (ruta2) -- node[etq, right=1pt]{no} (nc);

  \node[font=\diag, text=ESCOMGris, anchor=north west] at (0,-6.2)
    {Propósitos: sello del servidor, firma individual y sello de tiempo. El despliegue documentado fija una sola raíz con los tres.};
  \node[nota] at (2.45,-7.25) {Ruta no disponible: la validación sigue y solo cambia el estado actual};
  \node[nota] at (8.0,-7.25) {Nunca se siguen direcciones encontradas dentro del paquete};
\end{tikzpicture}

  \caption{Decisión sobre una raíz de confianza. Solo la lista fija de la versión de Qadra convierte una raíz en ancla; una raíz que no está en la lista fija ni la devuelve la ruta deja el paquete como no compatible, no como alterado.}
  \label{fig:cripto-validador-raices}
\end{figure}

\paragraph{Tres tipos de paquete.} El validador reconoce los dos paquetes que Qadra ya exportaba y uno nuevo: el sello del servidor, la firma individual en su versión 1 y la firma individual en su versión 2, que añade un manifiesto sellado (\cref{fig:cripto-manifiesto}). El tipo se decide por el conjunto de archivos que contiene, no por su orden. Un archivo desconocido, uno obligatorio ausente o un documento ambiguo hacen que el paquete se informe como no compatible: el validador no adivina qué quiso exportarse.

\paragraph{Manifiesto sellado de la versión 2.} El sello de tiempo del documento cubre solo la huella del PDF; no demuestra que la firma separada, el consentimiento o el material de confianza ya existieran en ese instante. El manifiesto cierra esa brecha. Es un archivo JSON que enumera las catorce piezas de evidencia del paquete, cada una con su función, su huella SHA-256 y su tamaño, y que identifica el caso, el documento, su versión, la operación de aceptación y al firmante. Inmediatamente después de verificar la firma y obtener el sello del PDF, dentro de la misma operación de aceptación, el servidor construye el manifiesto con los bytes finales y pide a la TSA un segundo sello sobre su huella. Así, un solo sello ata todas las piezas a un instante emitido por la TSA.

El manifiesto se escribe en la forma canónica de JSON definida por RFC 8785 \cite{rundgren2020}. JSON admite muchas escrituras equivalentes del mismo contenido (orden de claves, espacios, escapes); la forma canónica fija una sola, de modo que cualquiera puede recalcular la huella sobre los mismos bytes y el validador rechaza un manifiesto que no esté exactamente en esa forma. El paquete agrega además la lista de revocación del emisor de la TSA, capturada al sellar, para comprobar que la TSA no estaba revocada, y un recibo de aceptación que repite la huella del manifiesto y el instante registrado. El manifiesto no se incluye a sí mismo, ni a su sello, ni al recibo, ni al instructivo, ni al registro de tiempos de Qadra, cuyo instante de aceptación se fija al confirmar la operación. Ese instante sigue siendo un registro de la aplicación, pero en la versión 2 queda acotado: el sello del documento no puede ser posterior al del manifiesto, ni este a la aceptación. El perfil admite un solo firmante; varios firmantes quedan fuera de esta entrega.

\begin{figure}[!htbp]
  \centering
  % Fuente: docs/adr/0085-in-browser-evidence-validation.md (perfiles y manifiesto);
% crates/application/src/evidence.rs, crates/application/src/evidence/individual.rs
% y crates/application/src/documents/signing_manifest.rs.
\begin{tikzpicture}[
  e/.style={cajaClara, text width=3.15cm, minimum height=0.62cm, inner sep=2pt},
  n/.style={cajaGris, text width=3.15cm, minimum height=0.62cm, inner sep=2pt},
  r/.style={font=\diag, text=ESCOMGris}]

  % Zonas
  \draw[zona] (0,0) rectangle (3.75,-7.55);
  \draw[zona] (4.05,0) rectangle (7.8,-7.55);
  \draw[zona] (8.1,0) rectangle (16,-7.55);
  \node[zonaLabel] at (0,0) {A · TODOS LOS PAQUETES};
  \node[zonaLabel] at (4.05,0) {B · FIRMA INDIVIDUAL};
  \node[zonaLabel] at (8.1,0) {C · SOLO FIRMA INDIVIDUAL, VERSIÓN 2};

  % A: comunes
  \foreach \y/\st/\t/\d in {
      -0.95/e/{\texttt{<documento>}}/{documento firmado},
      -1.75/e/{\texttt{<documento>.sig}}/{firma separada},
      -2.55/e/{\texttt{<documento>.tsr}}/{sello de tiempo del documento},
      -3.35/e/{\texttt{certificado.pem}}/{certificado firmante},
      -4.15/e/{\texttt{ca.pem}}/{raíz emisora},
      -4.95/e/{\texttt{crl.pem}}/{CRL de esa raíz},
      -5.75/e/{\texttt{tsa-chain.pem}}/{raíz de la TSA},
      -6.75/n/{\texttt{INSTRUCCIONES.md}}/{guía para OpenSSL}} {
    \node[\st] at (1.875,\y) {\t\\[-1pt]{\diag \d}};
  }

  % B: firma individual
  \foreach \y/\st/\t/\d in {
      -1.1/e/{\texttt{certificate-binding.bin}\\[-1pt]\texttt{certificate-binding.sig}}/{vinculación y su firma},
      -2.25/e/{\texttt{binding-root.der}}/{raíz al registrar el vínculo},
      -3.05/e/{\texttt{binding-crl.der}}/{CRL al registrar el vínculo},
      -4.2/e/{\texttt{document-consent.bin}\\[-1pt]\texttt{document-consent.sig}}/{consentimiento y su firma},
      -6.75/n/{\texttt{individual-signing.json}}/{registro de tiempos de Qadra}} {
    \node[\st] at (5.925,\y) {\t\\[-1pt]{\diag \d}};
  }

  % C: manifiesto sellado
  \node[cajaOsc, text width=3.9cm, align=left] (man) at (10.45,-1.75)
    {\textbf{\texttt{manifiesto.json}}\\[1pt]
     {\diag JSON canónico (RFC 8785). Lista los 14 archivos en azul claro (dos cajas agrupan un archivo y su firma) con su función, SHA-256 y tamaño, y nombra caso, documento, versión, operación y firmante}};
  \node[caja, text width=2.55cm] (tsr) at (14.5,-1.75)
    {\textbf{\texttt{manifiesto.json.tsr}}\\ {\diag sello RFC 3161 sobre la huella del manifiesto}};
  \node[e, text width=2.55cm] (tcrl) at (14.5,-3.95)
    {\texttt{tsa-crl.pem}\\[-1pt] {\diag CRL del emisor de la TSA}};
  \node[n, text width=3.9cm, align=left] (rec) at (10.45,-4.05)
    {\texttt{recibo-aceptacion.json}\\[-1pt] {\diag operación, huella del manifiesto e instante de aceptación; se genera al exportar y no lleva sello}};

  \draw[flecha] (tsr) -- node[etq, above=1pt]{sella} (man);
  \draw[flecha] (tcrl) -- node[etq, left=1pt, align=right]{TSA no\\ revocada} (tsr);
  \draw[flecha] (rec) -- node[etq, right=1pt]{cita su huella} (man);

  % Orden de los instantes
  \node[r, anchor=west] at (8.35,-5.15) {Orden que exige la versión 2:};
  \node[caja, text width=1.8cm] (t1) at (9.3,-6.0) {sello del documento};
  \node[caja, text width=1.8cm] (t2) at (12.05,-6.0) {sello del manifiesto};
  \node[cajaOsc, text width=1.8cm] (t3) at (14.8,-6.0) {aceptación registrada};
  \node at ($(t1.east)!0.5!(t2.west)$) {$\le$};
  \node at ($(t2.east)!0.5!(t3.west)$) {$\le$};

  % Leyenda
  \node[e, text width=0.4cm, minimum height=0.3cm] (l1) at (8.6,-7.05) {};
  \node[r, anchor=west] at (8.9,-7.05) {cubierta por el manifiesto};
  \node[n, text width=0.4cm, minimum height=0.3cm] (l2) at (11.9,-7.05) {};
  \node[r, anchor=west] at (12.2,-7.05) {auxiliar o registro, sin cubrir};
\end{tikzpicture}

  \caption{Contenido de los tres tipos de paquete y cobertura del manifiesto. El sello del servidor contiene A, con \texttt{tsa-chain.pem} opcional; la firma individual, versión 1, contiene A y B; la versión 2 añade C. El manifiesto lista los catorce archivos en azul claro, dos de ellos agrupados con su firma, y su propio sello lo ata a un instante de la TSA.}
  \label{fig:cripto-manifiesto}
\end{figure}

\paragraph{Cómo se expresa el resultado.} El informe se divide en ocho componentes, desde la forma del paquete hasta las fechas y registros, y cada uno obtiene uno de cinco resultados: correcto, fallido, falta evidencia para decidirlo, no compatible o no aplica (\cref{fig:cripto-validador-resultados}). Todos los componentes son obligatorios, salvo dos que se muestran como \emph{No aplica}: la vinculación y el consentimiento en el sello del servidor, y el manifiesto fuera de la versión 2. El resultado global toma el primero que aplica: \emph{Fallido} si algún componente obligatorio falló; \emph{No compatible} si alguno no es compatible; \emph{Parcial} si falta evidencia en alguno; y \emph{Completo} en otro caso. Un estado actual, con la vigencia y la revocación al día de hoy según el reloj del equipo, se presenta aparte como informativo y nunca cambia el resultado.

\begin{figure}[!htbp]
  \centering
  % Fuente: docs/adr/0085-in-browser-evidence-validation.md (resultados por componente);
% crates/evidence-verify/src/report.rs y web/apps/qadra/src/lib/evidence/presentation.mjs.
\begin{tikzpicture}[
  comp/.style={caja, text width=3.75cm, minimum height=0.6cm, align=left},
  esc/.style={caja, text width=3.3cm, minimum height=0.95cm, align=left},
  r/.style={font=\diag, text=ESCOMGris}]

  % Componentes
  \draw[zona] (0,0) rectangle (4.45,-7.7);
  \node[zonaLabel] at (0,0) {COMPONENTES};
  \foreach \i/\t in {
      1/{Paquete},
      2/{Sello de tiempo del documento},
      3/{Firma del documento},
      4/{Vinculación y consentimiento\\{\diag no aplica al sello del servidor}},
      5/{Certificados y raíz de confianza},
      6/{Revocación},
      7/{Manifiesto y su sello de tiempo\\{\diag solo en la versión 2}},
      8/{Fechas y registros}} {
    \node[comp] (c\i) at (2.225,{-0.2-0.88*\i}) {\t};
    \draw[ESCOMAzulOscuro, line width=0.6pt] (c\i.east) -- (4.8,{-0.2-0.88*\i});
  }
  \draw[ESCOMAzulOscuro, line width=0.6pt] (4.8,-1.08) -- (4.8,-7.24);

  % Cinco resultados
  \node[cajaClara, text width=2.6cm, align=left, minimum height=3.3cm] (cinco) at (6.55,-4.16)
    {\textbf{Cada uno obtiene uno de cinco resultados}\\[3pt]
     Correcto\\ Fallido\\ Falta evidencia para decidirlo\\ No compatible\\ No aplica};
  \draw[flecha] (4.8,-4.16) -- (cinco.west);

  % Precedencia
  \node[r, anchor=south west, align=left] at (8.4,-0.5) {Resultado global: se toma el primero que aplica};
  \node[esc, cajaError] (f) at (10.15,-1.2) {\textbf{Fallido}: algún componente obligatorio falló};
  \node[esc] (nc) at (10.15,-2.7) {\textbf{No compatible}: un obligatorio no es compatible o ningún perfil coincide};
  \node[esc] (p) at (10.15,-4.2) {\textbf{Parcial}: falta evidencia en algún obligatorio};
  \node[esc, cajaOsc] (c) at (10.15,-5.7) {\textbf{Completo}: todos los obligatorios comprobados};
  \draw[flecha] (cinco.east) -- ++(0.2,0) |- (f.west);
  \draw[flecha] (f) -- node[etq, right=1pt]{si no} (nc);
  \draw[flecha] (nc) -- node[etq, right=1pt]{si no} (p);
  \draw[flecha] (p) -- node[etq, right=1pt]{si no} (c);

  % Informativo y aclaración
  \node[cajaGris, text width=2.4cm] (act) at (14.65,-1.2) {\textbf{Estado actual}\\ vigencia y revocación hoy, con el reloj del equipo};
  \draw[flechaSuave] (act.west) -- node[etq, below=3pt, align=center]{informativo:\\ no cambia\\ el resultado} (f.east);
  \node[caja, text width=2.4cm, align=left] (nota) at (14.65,-4.7)
    {\textbf{Firma correcta con perfil de certificado fallido}: el resultado no es completo, pero eso no significa por sí solo que el PDF se haya alterado};

  \node[r, text width=10.6cm, align=left, anchor=north west] at (5.25,-6.65)
    {Se separan tres preguntas: si la firma corresponde a los bytes exactos, si los certificados cumplen el perfil y encadenan a una raíz fija, y si falta evidencia para decidir.};
\end{tikzpicture}

  \caption{Modelo de resultados. Cada componente se evalúa por separado y el resultado global sigue una precedencia fija; el estado actual es informativo y no interviene.}
  \label{fig:cripto-validador-resultados}
\end{figure}

El informe separa tres preguntas que una sola palabra confundiría: si la firma corresponde a los bytes exactos del documento, si los certificados cumplen el perfil estricto y encadenan a una raíz fija, y si falta evidencia para decidir. Por eso, un certificado que no cumple el perfil estricto produce un resultado que no es completo, pero no significa que el PDF se haya alterado; el panel lo aclara cuando la firma es correcta. Cada comprobación se evalúa en su propio instante: el certificado de la TSA, en la fecha de cada sello; el vínculo, al registrarse; la firma y el consentimiento del titular, al aceptarse. Así, una revocación posterior no borra la validez histórica y se muestra solo en el estado actual. Los paquetes históricos se leen sin reescribirse: si no traen la lista de revocación del emisor de la TSA y la lista general no cubre la fecha del sello, quedan en \emph{Parcial}. Como el validador aplica perfiles estrictos a todos los certificados y revisa la TSA en la fecha de cada sello, a diferencia del verificador del servidor, su veredicto sobre un paquete histórico puede diferir del que mostró el servidor. Las decisiones constan en el ADR~0085 y la implementación, en la \cref{sec:impl-validador}.

\subsection{Cifrado en reposo}\label{sec:cifrado-reposo}

Los documentos se cifran con AES-256-GCM conforme a NIST SP 800-38D \cite{dworkin2007}, mediante envoltura de claves (\emph{envelope encryption}) conforme a NIST SP 800-38F \cite{dworkin2012}. Una clave maestra del servidor (KEK) envuelve una clave de datos (DEK) aleatoria y propia de cada versión (\cref{fig:cripto-sobre}). Los datos adicionales autenticados (AAD) del documento contienen su UUID y su versión, así que descifrarlo bajo otro contexto falla. El cifrado no detecta por sí solo que se restaure un respaldo antiguo, pero válido, de documentos y auditoría; esa protección requiere un anclaje externo, que queda pendiente.

Cada operación usa un nonce aleatorio de 96 bits generado con un generador aleatorio criptográfico (CSPRNG), en lugar de un contador. Un contador exige un estado persistente; si se reinicia por un error, una restauración de respaldo o una condición de carrera, repite un valor y GCM pierde tanto la confidencialidad como la autenticación. Con nonces aleatorios, SP 800-38D limita cada clave a $2^{32}$ cifrados. Como cada DEK cifra una sola versión, ese límite solo afecta a la KEK, que envuelve todas las DEK y cifra algunos secretos pequeños (TOTP e invitaciones); por eso debe considerarse al planear su rotación.

Se descartaron dos alternativas. Una clave única para todo el repositorio concentra el riesgo y obliga a recifrar todo al rotarla. Derivar una clave por documento con HKDF evita guardar claves envueltas, pero el compromiso de la clave maestra expone todas las derivadas y la unicidad del contexto reintroduce un problema de estado. Con la envoltura, el compromiso de una DEK expone una sola versión y el de la KEK exige además acceso al almacenamiento cifrado; es el patrón de AWS KMS, Google Cloud KMS y Azure Key Vault. El prototipo implementa y documenta la rotación de la KEK; ejecutarla periódicamente queda fuera del alcance de la demostración.

\begin{figure}[!htbp]
  \centering
  \begin{tikzpicture}[
  dato/.style={caja, text width=3.0cm, minimum height=0.85cm},
  op/.style={cajaClara, text width=3.3cm, minimum height=1.1cm},
  sec/.style={cajaOsc, text width=3.0cm, minimum height=0.85cm},
  campo/.append style={font=\scriptsize\diagsans},
  tam/.append style={font=\diag}]

  % Zonas
  \draw[zona] (0,0) rectangle (10.6,-4.4);
  \draw[zona] (11.1,0) rectangle (16,-4.4);
  \node[zonaLabel] at (0,0) {SERVIDOR DE LA API};
  \node[zonaLabel] at (11.1,0) {POSTGRESQL};

  % Entradas
  \node[dato] (doc) at (1.85,-1.0) {Versión del documento en claro};
  \node[sec] (dek) at (1.85,-2.25) {DEK aleatoria de 32~B, una por versión};
  \node[sec] (kek) at (1.85,-3.5) {KEK del servidor (configuración)};

  % Operaciones
  \node[op] (g1) at (6.6,-1.0) {\textbf{AES-256-GCM}\\ nonce nuevo de 12 B\\ AAD = UUID + versión};
  \node[op] (g2) at (6.6,-3.5) {\textbf{AES-256-GCM}\\ nonce nuevo de 12 B\\ AAD vacío};
  \draw[flecha] (doc) -- node[etq, above]{texto claro} (g1);
  \draw[flecha] ($(dek.east)+(0,0.2)$) -| node[etq, above, pos=0.3]{clave} (g1.south);
  \draw[flecha] ($(dek.east)+(0,-0.2)$) -| node[etq, below, pos=0.3]{texto claro} (g2.north);
  \draw[flecha] (kek) -- node[etq, above]{clave} (g2);

  % Salidas
  \node[dato] (c1) at (13.55,-1.0) {Documento cifrado};
  \node[dato] (c2) at (13.55,-3.5) {DEK envuelta};
  \draw[flecha] (g1) -- (c1);
  \draw[flecha] (g2) -- (c2);

  % Formato común
  \node[anchor=south west, font=\diag\bfseries, text=ESCOMAzulOscuro, inner sep=1pt]
    at (0,-5.05) {Formato de ambos paquetes cifrados};
  \Tira{0}{-6.15}{2.0/campo/{nonce}//12 B,
                  4.6/campoRel/{texto cifrado}//,
                  2.4/campoFirma/{etiqueta GCM}//16 B}

  % Rotación
  \node[paso, font=\scriptsize\diagsans, text width=6.1cm, anchor=north west] at (9.75,-5.0)
    {\textbf{Rotación de la KEK.} Cada DEK se abre con la KEK vigente y se vuelve a envolver con la nueva. El documento cifrado no cambia, así que el costo no depende de su tamaño.};
\end{tikzpicture}

  \caption{Envoltura de claves. Cada versión tiene su DEK; para los documentos, la KEK solo envuelve claves de 32 B, por lo que rotarla no toca los documentos cifrados.}
  \label{fig:cripto-sobre}
\end{figure}

\subsection{Autenticación y control de acceso}\label{sec:autenticacion}

Las contraseñas se almacenan con Argon2id (RFC 9106) \cite{biryukov2015,biryukov2021}, la variante recomendada por OWASP \cite{owasp_pwstorage2024} y NIST SP 800-63B Rev 4 \cite{nist2024}. Los hashes nuevos usan $m = 256$~MiB, $t = 3$, $p = 1$, sal de 16 bytes y salida de 32 bytes. El objetivo es que un hash nuevo cueste entre 0.5 y 1.0 segundos en el hardware de referencia; con tres iteraciones tardó unos 0.66~s en el equipo de referencia, dentro de esa banda (\cref{sec:calibracion-credenciales}). La banda no cubre el inicio de sesión completo ni su comportamiento bajo carga, y el hardware de un despliegue debe medirse antes de usarlo. Los hashes anteriores conservan sus parámetros y no se reescriben al verificarlos. Argon2id se prefirió a bcrypt por su dureza de memoria: bcrypt usa unos 4 KiB por intento, mientras que reservar 256 MiB encarece los ataques masivos con GPU o ASIC. Argon2, ganador de la Password Hashing Competition de 2015, tiene la variante Argon2id, que combina Argon2i, resistente a canales laterales, con Argon2d, resistente a hardware especializado.

El acceso combina correo y contraseña con un segundo factor TOTP (RFC 6238) \cite{mraihi2011} o un código de recuperación de un solo uso; el recorrido completo aparece en la \cref{fig:cu-seq-sesion}. El TOTP usa HMAC-SHA1, seis dígitos, periodos de 30 segundos, tolerancia de un periodo a cada lado y secreto de 20 bytes. Para comprobar un código, la API abre el secreto cifrado con la KEK. SHA-1 está desaconsejado para firmas por sus colisiones, pero la seguridad de HMAC no depende de esa propiedad; NIST SP 800-107 \cite{dang2012} lo mantiene aceptable para autenticación, y es el único algoritmo que todas las aplicaciones autenticadoras soportan. Cada intento del segundo factor consume el desafío, y una reclamación atómica en Redis impide reutilizar un código aceptado. El acceso por certificado del Owner, deshabilitado por defecto, conserva límites propios y no sustituye al segundo factor. Ambos recorridos emiten un token opaco y revocable; los certificados de acceso del Owner y de firma documental mantienen propósitos separados y no acreditan identidad civil.

El alta inicial crea al primer Owner solo mientras no existan usuarios; las altas posteriores requieren un Owner autenticado. En la creación directa, la API devuelve una sola vez el secreto TOTP, su URI \texttt{otpauth} y ocho códigos de recuperación. El secreto se guarda cifrado con AES-256-GCM y ligado al UUID del usuario, y los códigos con Argon2id. El alta por invitación separa preparación y confirmación: el destinatario elige su contraseña, recibe el material MFA y la cuenta se crea solo al confirmar un código TOTP de su dispositivo (\cref{sec:impl-invitaciones}). Qadra muestra la URI para configurarla a mano; presentarla como código QR queda como mejora de la interfaz.

\begin{table}[!htbp]
  \centering
  \caption{Permisos por rol sobre administración, documentos, etapas y directorio.}
  \label{tab:permisos-implementados}
  \footnotesize
  \begin{tabularx}{\linewidth}{@{}X C{0.14\linewidth} C{0.16\linewidth} C{0.16\linewidth}@{}}
    \toprule
    \textbf{Operación} & \textbf{Owner} & \textbf{Litigante} & \textbf{Paralegal} \\
    \midrule
    Consultar administración e historial & Sí & Sí & Sí \\
    Registrar o administrar expediente & Sí & Sí & No \\
    Consultar metadatos documentales & Sí & Sí & Sí \\
    Cargar o añadir versión & Sí & Sí & Sí \\
    Clasificar documento & Sí & Sí & Sí \\
    Sellar documento & Sí & Sí & No \\
    Verificar documento & Sí & Sí & Sí \\
    Exportar evidencia & Sí & Sí & Sí \\
    Consultar etapa e historial & Sí & Sí & Sí \\
    Adoptar o avanzar etapa & Sí & Sí & No \\
    Consultar directorio e historial & Sí & Sí & Sí \\
    Crear, editar o archivar ficha & Sí & Sí & No \\
    Verificar auditoría completa & Sí & No & No \\
    \bottomrule
  \end{tabularx}

  \smallskip
  \parbox{\linewidth}{\footnotesize Cliente no tiene ninguna de estas operaciones; solo consulta un resumen de sus expedientes.}
\end{table}

La autorización usa el rol y el estado vigentes de cada cuenta. Owner ve todos los expedientes; Litigante y Paralegal, solo los asignados; Cliente, solo un resumen básico de los suyos. Owner y Litigante administran perfiles y estado; las asignaciones quedan reservadas a Owner. En documentos, Litigante y Paralegal necesitan asignación vigente además del permiso de la \cref{tab:permisos-implementados}; Cliente queda denegado aunque esté asignado, y la asociación exacta entre documento y expediente se comprueba también para Owner. Una revocación surte efecto en la siguiente operación documental que aún no se haya confirmado.

Estos permisos se combinan con pertenencia vigente, recurso exacto y estado administrativo. Un expediente cerrado conserva las lecturas y la evidencia, pero rechaza nuevas cargas, versiones, clasificaciones, sellos, ediciones del perfil, cambios del directorio y registros de etapa. Reactivarlo no modifica etapas ni reactiva participantes archivados.

\subsection{Postura legal del prototipo}\label{sec:postura-legal}

La firma del prototipo es una firma electrónica simple conforme al artículo 89 del Código de Comercio \cite{ccom2018}, porque sus certificados no provienen de un PSC autorizado (artículo 97, requisito 5). Es admisible como prueba, ya que el artículo 89 no condiciona la admisibilidad a un PSC, y la equivalencia funcional del artículo 100 permite argumentar la fiabilidad del método con estándares internacionales: RSA-3072 según NIST SP 800-57 \cite{barker2020} y SHA-256 según FIPS 180-4 \cite{nist2015fips1804}. La diferencia con una firma avanzada está en la carga de la prueba: quien invoca la firma simple debe acreditar la fiabilidad del método, mientras que la avanzada goza de presunción legal. Del mismo modo, el token de la TSA local demuestra el mecanismo RFC 3161, pero no aporta la fecha cierta ni la independencia de tercero que daría un sello de un PSC autorizado conforme a la NOM-151. Los puertos criptográficos permiten integrar certificados y autoridades externas mediante adaptadores, pero adoptarlos exige evaluar identidad, custodia de claves, políticas y confianza en el proveedor; cambiar el certificado no produce por sí solo una categoría jurídica distinta.
